\section{Results}
\label{sec:results}

% 为同时评估评分性能与格式遵循，本文的主要分析均采用严格解析，
%仅接受符合预定义输出格式的分数；可恢复其他可解释输出的宽松解析结果见附录 A。

All primary analyses use strict parsing to evaluate both scoring performance and format 
compliance, accepting only schema-conforming scores; relaxed-parsing results are 
reported in Appendix A.

% \paragraph{RQ1：SO 与 RS 在评分性能和 rationale 建模方面有何差异？}
\paragraph{RQ1: How Do SO and RS Differ in Scoring Performance and Rationale Modeling?}

% Figure~2 展示了四个序数任务上的 QWK。采用 score-only training
% 时，direct-score 和 rationale-first inference 的平均 QWK 分别为
% 0.747 和 0.712；采用 rationale-and-score training 时，分别为
% 0.696 和 0.679。所有任务均呈现相同趋势：score-only training
% 优于 rationale-and-score training，direct-score inference 优于
% rationale-first inference。

Figure~2 reports QWK on the four ordinal tasks. With score-only training, 
direct-score and rationale-first inference achieve mean QWK scores of 0.747 and 0.712, 
respectively. With rationale-and-score training, the corresponding scores are 0.696 and 0.679. 
This pattern is consistent across tasks: score-only training outperforms rationale-and-score 
training, and direct-score inference outperforms rationale-first inference.

% Figure~3 展示了二分类任务上的 Macro-F1。采用 score-only training
% 时，direct-score 和 rationale-first inference 的平均结果接近，
% 分别为 0.928 和 0.923。采用 rationale-and-score training 时，
% rationale-first inference 的平均 Macro-F1 为 0.922，明显高于
% direct-score inference 的 0.843。因此，二分类任务中的训练效果
% 取决于推理接口：score-only training 仅在 direct-score inference
% 下表现出明显优势。

Figure~3 reports Macro-F1 on the binary tasks. With score-only training, direct-score and 
rationale-first inference perform similarly, scoring 0.928 and 0.923, respectively. 
With rationale-and-score training, rationale-first inference reaches 0.922, compared 
with 0.843 under direct-score inference. The effect of the training objective therefore 
depends on the inference protocol, with a clear score-only advantage appearing only under 
direct-score inference.

% \textbf{Rationale supervision 增强了教师 rationale 拟合，但这一优势并未转化为更好的评分性能。}
% 我们通过 teacher-forced rationale NLL 衡量模型对教师 rationale
% 的建模程度。未经任务微调的 Qwen3-4B（Base）平均 NLL 为
% 3.4480，经过 score-only 和 rationale-and-score training 后
% 分别降至 1.6964 和 0.7793。Rationale-and-score training 在
% 所有任务上均取得最低 NLL，表明直接监督 rationale 能够显著
% 增强对教师 rationale token 的预测，但不能保证更好的评分结果。

\textbf{Rationale supervision improves teacher-rationale fitting, but this 
advantage does not translate into better scoring.} 
We measure teacher-rationale modeling using teacher-forced rationale NLL. 
The original Qwen3-4B without task-specific fine-tuning (Base) obtains a mean NLL of 3.4480, 
which decreases to 1.6964 with score-only training and 0.7793 
with rationale-and-score training. 
Rationale-and-score training achieves the lowest NLL on every task, 
showing that direct rationale supervision improves teacher-token prediction 
but does not guarantee better scoring performance.

% 各任务的 teacher-forced rationale NLL。Base 表示未经任务微调的
% Qwen3-4B，数值越低表示对教师 rationale token 的拟合越好。

\begin{table}[t]
\centering
\scriptsize
\setlength{\tabcolsep}{3pt}
\resizebox{\columnwidth}{!}{%
\begin{tabular}{lrrrr}
\toprule
Task & Base & SO & RS & RS $-$ SO \\
\midrule
Actionability         & 3.3984 & 1.5891 & \textbf{0.7894} & $-$0.7997 \\
Grounding Specificity & 3.6439 & 1.6757 & \textbf{0.6761} & $-$0.9996 \\
Helpfulness           & 3.5397 & 1.7811 & \textbf{0.8856} & $-$0.8955 \\
Verifiability         & 3.5677 & 1.9243 & \textbf{0.8401} & $-$1.0842 \\
Coherence             & 3.4272 & 1.4634 & \textbf{0.7852} & $-$0.6782 \\
Positioning Check     & 3.1925 & 1.5362 & \textbf{0.7112} & $-$0.8250 \\
Positioning Type      & 3.3664 & 1.9051 & \textbf{0.7676} & $-$1.1375 \\
\midrule
Average               & 3.4480 & 1.6964 & \textbf{0.7793} & $-$0.9171 \\
\bottomrule
\end{tabular}%
}

% 各任务的 teacher-forced rationale NLL；数值越低越好。

\caption{Teacher-forced rationale NLL by task; lower is better.}
\label{tab:rationale-nll}
\end{table}

% \textbf{Rationale-and-score training 生成了质量更高的 rationale。}
% 外部 LLM 比较了 350 组配对 rationale，其中 248 组（70.9\%）
% 更偏好 rationale-and-score training 的结果，25 组（7.1\%）
% 更偏好 score-only training，其余 77 组（22.0\%）基本相当。
% 这一偏好在所有评测任务上均保持一致。Rationale-and-score
% training 在评分标准理解、证据支撑、关键信息覆盖和无依据陈述
% 控制四个维度上均取得更高分，其中关键信息覆盖的差距最大。
% 其四维平均分为 2.832，高于 score-only training 的 2.400；
% 与预测分数的一致率也更高，分别为 96.6\% 和 74.3\%。

\textbf{Rationale-and-score training produces higher-quality rationales.} 
External LLMs compared 350 paired rationales and preferred those from rationale-and-score 
training in 248 cases (70.9\%). Score-only rationales were preferred in 25 cases (7.1\%), 
while 77 cases (22.0\%) were comparable. This preference was consistent across all 
evaluation tasks. Rationale-and-score training scored higher on rubric understanding, 
evidence grounding, key-information coverage, and unsupported-claim control, 
with the largest difference in key-information coverage. Its mean score across the four 
dimensions was 2.832, compared with 2.400 for score-only training, and its rationale--score 
consistency was also higher at 96.6\% versus 74.3\%.

% 外部 LLM 对 rationale 的整体偏好、质量评分和分数一致性评估。
% 两条实线分别分隔 Pairwise Preference、Quality 和 Consistency。

\begin{table*}[t]
\centering
\small
\setlength{\tabcolsep}{9pt}
\renewcommand{\arraystretch}{1.05}
\begin{tabular}{llrrr}
\toprule
Analysis & Measure & SO & RS & Comparison \\
\midrule
Pairwise Preference
& Preferred rationale, $n$ (\%)
& 25 (7.1) & \textbf{248 (70.9)} & Tie: 77 (22.0) \\
\midrule
Quality
& Rubric understanding
& 2.379 & \textbf{2.848} & +0.469 \\
& Evidence grounding
& 2.495 & \textbf{2.837} & +0.342 \\
& Key-information coverage
& 2.273 & \textbf{2.833} & +0.560 \\
& Unsupported-claim control
& 2.453 & \textbf{2.811} & +0.358 \\
& Four-dimension mean
& 2.400 & \textbf{2.832} & +0.433 \\
\midrule
Consistency
& Supported (\%)
& 74.3 & \textbf{96.6} & +22.3 pp \\
& Partially supported (\%)
& 20.0 & \textbf{3.1} & $-$16.9 pp \\
& Not supported (\%)
& 5.7 & \textbf{0.3} & $-$5.4 pp \\
\bottomrule
\end{tabular}

% 外部 LLM 对 350 组 rationale 的配对评估结果。质量维度采用
% 1--3 分量表。Comparison 在偏好行报告平局数量，其余行报告
% RS 减 SO 的差值，其中比例差以百分点表示。

\caption{External-LLM evaluation of 350 paired rationales generated
under rationale-first inference. Quality dimensions use a 1--3 scale.
For pairwise preference, Comparison reports tied judgments; elsewhere,
it reports RS minus SO, with percentage differences in percentage points.}
\label{tab:rationale-quality-consistency}
\end{table*}

% \subsection{RQ2：为什么 RF 在序数任务上弱于 DS？}
\paragraph{RQ2: Why Does RF Underperform DS on Ordinal Tasks?}

% 我们在四个序数任务和三个 seed 上追踪 DS 到 RF 的样本级预测
% 转移，每种训练目标包含 11,364 个 seed-level 配对事件。切换到
% RF 后，SO 和 RS 的平均 QWK 分别下降 0.035 和 0.017。在 SO
% 条件下，准确率下降 3.52 个百分点，严重错误增加 3.98 个百分点；
% 在 RS 条件下，准确率虽然提高 3.49 个百分点，严重错误仍增加
% 4.26 个百分点。由于 QWK 对跨档距离赋予更高惩罚，严重错误的
% 增加抵消了正确样本增加带来的收益。格式无效率变化不足
% 0.25 个百分点，并非 QWK 下降的主要来源。
We traced sample-level DS-to-RF transitions across four ordinal tasks
and three seeds, yielding 11,364 paired events per supervision
objective. Switching to RF reduced mean QWK by 0.035 under SO and
0.017 under RS. Under SO, accuracy decreased by 3.52 percentage points,
while severe errors increased by 3.98 points. Under RS, accuracy rose
by 3.49 points, yet severe errors increased by 4.26 points. Because QWK
penalizes cross-level errors more heavily, this increase offset the
gain in exact predictions. Invalid-output rates changed by less than
0.25 points and were not the main source of the QWK decline.

% 为定位严重错误的来源，我们抽取了 296 个预测发生有害或有益
% 转移的案例，其中 278 个形成有效裁判共识。裁判认为 158 条
% RF rationale 支持错误分数，120 条支持 gold score。句子级分析
% 显示，评分标准误用最为常见，分数映射错误、证据误读和事实错误
% 也频繁出现。首个错误平均约出现在第 2 句，约 69\% 的案例在
% 前两句内已经出现错误。这些结果表明 rationale 方向与预测转移
% 高度相关，但仅凭观察性审计仍不能确定因果关系。
To localize these severe errors, we audited 296 harmful or beneficial
transition cases, of which 278 reached valid judge consensus. The
judges found that 158 RF rationales supported an incorrect score,
whereas 120 supported the gold score. Rubric misapplication was the
most frequent sentence-level error, followed by score-mapping errors,
evidence misreadings, and factual errors. The first error appeared near
the second sentence on average, and about 69\% occurred within the
first two sentences. These observations linked rationale direction to
prediction transitions, but did not establish causality.

% 因此，我们对 158 个支持错误分数的案例实施固定前缀干预。
% 条件 A 使用自然 RF 输出，条件 B 固定原错误 rationale 并重新
% 生成分数，条件 C 则固定修正后的 rationale。条件 B 在 157/158
% 个案例中复现了条件 A 的原预测，说明固定前缀本身很少改变分数。
% 从 B 切换到 C 后，158 个分数全部变化，其中 142 个（89.87\%）
% 恢复为 gold score，其余 16 个均缩小为相邻错误；MAE 也从
% 2.044 降至 0.101。该干预证明了所选案例的分数会响应显式
% rationale 内容，但不证明模型隐藏计算过程的完全忠实性。
We therefore applied a fixed-prefix intervention to the 158 cases
whose rationales supported an incorrect score. Condition A used the
natural RF output. Condition B fixed the original erroneous rationale,
whereas Condition C fixed its corrected counterpart before regenerating
the score. Condition B reproduced 157 of 158 original predictions,
showing that prefix fixation rarely changed the score by itself.
Replacing B with C changed every prediction and restored 142 cases
(89.87\%) to gold. The remaining 16 became adjacent errors, while MAE
fell from 2.044 to 0.101. This intervention establishes sensitivity to
the exposed rationale on the selected cases, not complete faithfulness
of the model's hidden computations.

% \subsection{RQ3：如何同时保持 rationale 质量与评分性能？}
\subsection{RQ3: How Can RS Preserve Rationale Quality and Scoring Performance?}

% ReaSon 首先针对 RF 路径中的 rationale 来源进行调整。在 RF 下，
% 它将四个序数任务的平均 QWK 从 0.679 提高到 0.707，并将 MAE
% 从 0.583 降至 0.534；分类 Macro-F1 基本不变，七任务平均从
% 0.783 提高到 0.799。在原 RS 的 338 个有害转移事件中，ReaSon
% 使 172 个（50.9\%）案例的误差减小。然而，其 DS 宽松解析
% 七任务平均仅变化 $-$0.001，严格格式无效率却达到 47.74\%。
% 因此，ReaSon 的收益主要集中在 RF，且未解决 DS 接口覆盖问题。
ReaSon first changed the rationale source used along the RF path.
It increased mean ordinal QWK from 0.679 to 0.707 and reduced MAE
from 0.583 to 0.534. Classification Macro-F1 remained nearly
unchanged, raising the seven-task mean from 0.783 to 0.799. ReaSon
also reduced the error in 172 of 338 original RS harmful transitions
(50.9\%). However, its relaxed DS mean changed by only $-$0.001,
while its strict invalid-output rate reached 47.74\%. The gains were
therefore concentrated under RF and did not resolve DS interface
coverage.

% 仅在单条 RF 序列内重新加权同样无法解决这一问题：在已完成的
% 34 个评测中，DS 产生了 2,387 个严格格式错误，其中 2,378 个
% 虽可由宽松解析恢复，仍表明模型没有稳定遵循 DS 协议。PaIR
% 通过配对 DS 和 RF 训练视图处理接口覆盖问题。其 DS 七任务平均
% 达到 0.833，高于 SO\_DS 的 0.825 和 RS\_DS 的 0.768；
% RF 平均为 0.797，高于 RS\_RF 的 0.783，但仍低于 SO\_RF
% 的 0.803。与使用相同配对数据的 Mix 相比，PaIR 在 DS 和 RF
% 下分别提高 0.013 和 0.006，表明分支级归一化带来了额外收益。
Reweighting a single RF sequence did not solve this problem. Across
34 completed evaluations, DS produced 2,387 strict-format failures,
although relaxed parsing recovered 2,378 of them. PaIR instead paired
DS and RF training views to provide explicit interface coverage.
Its seven-task DS mean reached 0.833, exceeding SO\_DS (0.825) and
RS\_DS (0.768). Its RF mean reached 0.797, exceeding RS\_RF (0.783)
but remaining below SO\_RF (0.803). Against Mix, which used the same
paired data, PaIR gained 0.013 under DS and 0.006 under RF. This matched
ablation supports an additional benefit from branch-wise normalization.

% 最后，RePaIR 联合 ReaSon 的 rationale 重生成与 PaIR 的配对接口
% 监督。RePaIR\_DS 的七任务平均为 0.832，比 SO\_DS 高 0.007，
% 仅比 PaIR\_DS 低 0.001；RePaIR\_RF 达到 0.817，分别比
% SO\_RF 和 PaIR\_RF 高 0.014 和 0.020。两种接口的严格格式
% 有效率均为 100\%。因此，RePaIR 基本保留了 PaIR 的 DS 收益，
% 同时补充了 ReaSon 在 RF 路径上的增益，在两个接口上取得了
% 最均衡的整体结果。现有结果验证了评分和格式收益，但 RePaIR
% 是否保留 ReaSon 的 rationale 质量收益仍需同样本审计。
Finally, RePaIR combined ReaSon's rationale regeneration with PaIR's
paired-interface supervision. RePaIR\_DS achieved a seven-task mean
of 0.832, which was 0.007 above SO\_DS and 0.001 below PaIR\_DS.
RePaIR\_RF reached 0.817, exceeding SO\_RF by 0.014 and PaIR\_RF
by 0.020. Both interfaces maintained 100\% strict-format validity.
RePaIR therefore preserved nearly all of PaIR's DS gain while adding
the RF benefit of rationale regeneration, yielding the most balanced
aggregate performance across the two interfaces. These results verify
the scoring and formatting gains. A matched audit is still needed to
test whether RePaIR retains ReaSon's rationale-quality improvement.
